\SourceLeaf{304}
\Needspace{6\baselineskip}\begin{center}\textsc{Gauss} an \textsc{Schumacher}. Göttingen, 15. Mai 1843.\end{center}

Ich füge meinem letzten Brief an Sie, mein teuerster Freund, noch ein paar Worte
hinzu, um einem Buch gerecht zu werden, das ich bisher, zum
Teil infolge eines vor zwanzig Jahren von Ihnen (freilich ganz unschuldig)
veranlassten Umstands, weniger beachtet habe, als es verdient.

Sie erinnern sich vielleicht noch, dass Sie mir damals in Ihrem Haus eine an sich
recht elegante geometrische Aufgabe von \textsc{Möbius} vorlegten, für die ich
fast augenblicklich eine Lösung niederschrieb, die nichts zu wünschen übrig ließ. Infolge
dieses Vorfalls hatte ich das Buch von \textsc{Möbius}, barycentrischer Calcul, in
dessen Vorrede jene Aufgabe ebenfalls erwähnt wurde, mit einer Art Vorurteil zur
Hand genommen. Ich zweifelte nämlich, ob es die Mühe wert sei, sich ein recht geschickt
erdachtes Rechenverfahren anzueignen, wenn sich damit nichts leisten
ließe, was sich ohne dieses Verfahren nicht ebenso leicht leisten lässt. Ich hatte deshalb das
Buch --- zumal ich es in einer Zeit erhielt, in der mich andere Dinge ganz
in Anspruch nahmen --- zunächst ohne große Erwartungen beiseitegelegt
und später, wie es mir mit manchem ergangen ist, völlig vergessen.

In den allerletzten Tagen (genauer gesagt gestern) fiel mir das Buch zufällig in die
Hände. Ich stellte bald mit großem Vergnügen fest, dass darin die Quintessenz der
Lehre von den Kegelschnitten auf knappstem Raum zusammengefasst ist und dass gerade sein
baryzentrischer Kalkül auf dem einfachsten Weg zur Lösung aller einschlägigen
Aufgaben führt. Beim Blättern fand ich darin auch manches andere Interessante,
zum Beispiel den schönen Satz von \textsc{Mac Laurin} und \textsc{Brackenridge},
der mir entweder unbekannt geblieben war oder den ich wieder vergessen hatte. Mit
seiner Hilfe lässt sich die Aufgabe, durch 5 gegebene Punkte einen Kegelschnitt zu legen, so
elegant lösen.
